\textbf{Elección y deducción del procedimiento.} La misma integración hidrostática da $p=p_b(T/T_b)^{g/(R\beta)}$. Después se usa la ecuación de estado con temperatura local, no la de la base.

\begin{enumerate}

\item Calcular la temperatura sin redondeo previo.

\[\fitmath{T=\mathrm{Tb}-\beta\cdot H}\]

\[\fitmath{T=(288.15)-(0.0065)\cdot (6000)\approx 249.15\,\mathrm{K}}\]

\item Calcular la presión de la capa.

\[\fitmath{p=\mathrm{pb}\cdot {\left(\frac{T}{\mathrm{Tb}}\right)}^{\frac{g}{R\cdot \beta}}}\]

\[\fitmath{p=(101325)\cdot {\left(\frac{(249.15)}{(288.15)}\right)}^{\frac{(9.80665)}{(287.05287)\cdot (0.0065)}}\approx 47181\,\mathrm{Pa}}\]

\item Dividir presión por constante y temperatura absoluta.

\[\fitmath{\rho=\frac{p}{R\cdot T}}\]

\[\fitmath{\rho=\frac{(47181.00219)}{(287.05287)\cdot (249.15)}\approx 0.6596968\,\mathrm{kg/m^{3}}}\]

\end{enumerate}
